\documentclass[12pt,a4paper]{article}
\usepackage[margin=24mm]{geometry}
\usepackage[T1]{fontenc}
\usepackage{lmodern,amsmath,amssymb,amsthm,booktabs,array,xcolor}
\usepackage[colorlinks=true,linkcolor=blue!55!black,citecolor=blue!55!black,urlcolor=blue!55!black]{hyperref}
\newtheorem{proposition}{Proposition}
\setlength{\parindent}{0pt}
\setlength{\parskip}{6pt}
\setlength{\emergencystretch}{2em}
\title{Whole-Contraction JIT Compilation\\for Numerical Hypervolume in Four Dimensions}
\author{Implementation, correctness, and performance report}
\date{4 October 2026}
\begin{document}
\maketitle
\begin{abstract}
We apply the compilation strategy of the older \texttt{HV4DMagnitude} prefix
solver to the newer numerical Chan implementation. Complete term contractions
run in Numba, and term states remain native throughout the Python geometric
recursion. The geometric cuts and compression schedule are retained.
Across five completed matched 4D cases, warm execution is 1.47--8.13 times
faster than the earlier hybrid backend, with median ratio 7.76. Both forty-point
cases now complete. The older compiled prefix solver remains substantially
faster on the larger spherical inputs: compilation removes interpreter and
conversion overhead, but not the new algorithm's intermediate-term growth.
We explain the finite-sum implementation, direct magnitude support, correctness
checks, timing method, and a separately measured 48.4-second fresh-cache start.
No new asymptotic bound is claimed.
\end{abstract}

\section{Purpose and scope}
The objective was to keep the new numerical algorithm understandable while
moving complete numerical loops into just-in-time (JIT) compiled code.
The motivating implementation in \texttt{HV4DMagnitude} compiles its complete
$2+2$ target-plane evaluation~\cite{old}. The earlier numerical Chan backend
compiled only short array helpers; term construction and merging remained
under Python control. Its Numba label therefore did not imply that the
expensive algebra was compiled.

The new 4D backend compiles that algebra while preserving the existing
recursive partition and scheduled compression~\cite{release}. It operates
in floating-point arithmetic and supports both ordinary hypervolume and
dominated-set $\ell_1$ magnitude. All end-to-end measurements in this report
are four-dimensional. The generic higher-dimensional backend remains available
as the hybrid reference; this report makes no performance claim for it.

Chan's orthant algorithm provides the theoretical
$O(n^{d/3}\operatorname{polylog}n)$ setting, which gives an exponent $4/3$
in four dimensions~\cite{chan}. The present work changes how numerical
operations execute. It neither improves that exponent nor establishes a new
global complexity proof for the older prefix implementation.

\section{The finite numerical problem}
For nonnegative upper corners $A\subseteq\mathbb R^4$, write
\[
 D(A)=\bigcup_{a\in A}\prod_{i=1}^4[0,a_i].
\]
The public interface receives these corners. Minimization points $p$ with
reference $r$ are converted to $a=r-p$; both conventions describe the same
ordinary hypervolume after reflection.

On each axis, sort the distinct coordinates
$0=c_0<c_1<\cdots<c_m$. Index zero represents the origin, and index $k>0$
represents the interval $(c_{k-1},c_k]$. Box membership is constant on each
product cell. Ordinary volume uses the masses
\[
 w(0)=0,\qquad w(k)=c_k-c_{k-1}\quad(k>0).
\]
For dominated-set magnitude in the anchored $\ell_1$ box setting~\cite{magnitude},
use instead
\[
 w(0)=1,\qquad w(k)=\tfrac12(c_k-c_{k-1})\quad(k>0).
\]
This is the product measure with an origin atom and half Lebesgue density on
each coordinate axis. The solver stores masses directly. Zero-coordinate
generators therefore remain relevant for magnitude even when their ordinary
four-dimensional volume is zero.

The state is a finite signed sum of terms. A term has the form
\begin{equation}
 T(x)=c\prod_{i\in V}w_i(x_i)
       \prod_{e\in E}\mathbf 1_{R_e}(x),
 \label{eq:term}
\end{equation}
where each pairwise predicate is an integer-cutoff condition such as
$x_j\le f(x_i)$ or $x_j\ge f(x_i)$. Cutoff arrays can be increasing or
decreasing. Both bound senses and both directions are retained; they cannot
in general be collapsed into one staircase per coordinate pair after
elimination. Coefficients and intermediate unary arrays may be signed.

In code, a term is simply a named tuple containing a floating coefficient,
a native dictionary of unary arrays, and a native dictionary of cutoff arrays.
A copied term receives new dictionaries but shares its arrays. Every mask or
cutoff update allocates a replacement array. This keeps the representation
sparse and prevents one recursive branch from modifying another.

\section{One compiled contraction}
Here a \emph{contraction} means summing out a variable from a finite product
expression. For a fixed assignment $z$ of the remaining variables, collect
all bounds on the variable $x$ to be removed. Its feasible interval is
\[
 L(z)=\max_j\ell_j(z),\qquad U(z)=\min_k u_k(z).
\]
Constant bounds account for the endpoints of the unary array. Monotonicity
allows a cutoff in the opposite orientation to be inverted by a binary search.
Each candidate bound depends on at most one remaining variable.

Define the numerical prefix array
\[
 P(t)=\sum_{0\le x<t}w(x),\qquad P(0)=0.
\]
If $L\le U$, the mass to insert into the remaining expression is exactly
\begin{equation}
 \sum_{x=L}^{U}w(x)=P(U+1)-P(L).
 \label{eq:prefix}
\end{equation}
If $L>U$, the contribution is zero. The code partitions assignments by the
first candidate attaining the maximum lower bound and the first candidate
attaining the minimum upper bound. Strict comparisons against earlier
candidates make these winner regions disjoint, including ties.

On a winner region, each prefix value is a unary function of a remaining
variable, or a constant. If the two prefix values depend on different variables,
their difference yields two signed terms. If they depend on the same variable,
one unary array holds the difference. The required comparisons become masks
or monotone pairwise constraints, so expression~\eqref{eq:term} is retained.
Consecutive eliminations implement the numerical pairwise recursion without
constructing a dense pair table.

\begin{proposition}[Finite-sum preservation]
In exact arithmetic, the contraction described above preserves the sum of
the represented expression. Repeated contractions, and block compression
formed from them, preserve that sum as well.
\end{proposition}
\begin{proof}
For each fixed assignment of the remaining variables, the first-winner rule
selects exactly one lower winner and one upper winner. If the resulting interval
is empty, the original sum and the emitted contribution are zero. Otherwise,
equation~\eqref{eq:prefix} is the telescoping definition of the prefix array.
Multiplying by the unchanged factors and distributing a difference preserves
the value. Summing over all assignments proves the first claim. Applying the
same equality repeatedly proves the second. The fine-cell blocks form a
disjoint partition of the current index interval on each axis. Imposing
inclusive block-membership bounds and then eliminating the fine indices
therefore preserves total mass when summed over the block indices.
\end{proof}

This proposition explains the algebra, not a rounding-error guarantee for
the float64 implementation. For a small magnitude example, let
$w=(1,\tfrac12,1)$, hence $P=(0,1,\tfrac32,\tfrac52)$. The interval
$[1,2]$ has mass $P(3)-P(1)=\tfrac32$, while $[0,2]$ has mass
$P(3)-P(0)=\tfrac52$. The origin atom is included exactly when index zero
is included. No integration routine is involved.

\subsection*{Normalization and merging}
Each emitted state is normalized by dividing every unary array by its first
nonzero entry and moving those factors into the coefficient. Identical states
can then be combined. The compiled implementation hashes states in canonical
key order and verifies equality by comparing the complete arrays; floating
arrays use bitwise comparison, as in the hybrid backend's byte-string keys.
Dictionary insertion order is irrelevant to this equality test, while the
candidate-bound order is preserved for tie breaking.

A lookup examines at most sixteen entries in its hash collision chain.
If no equal state is found, the new term is retained separately. A missed
merge does not discard a contribution, so this limits collision-search work
without introducing an approximation. It does not bound total term production
or the cost of reading an individual state's arrays. Cancelled terms are
compacted only after emission, keeping chain indices valid throughout merging.

\section{The recursion stays clear}
The geometric driver maintains a cell, the remaining hard orthants, and the
numerical state. First discard boxes disjoint from the current cell. For a
remaining box, coordinate $i$ is active when its upper index lies below the
cell's upper index. A box with no active coordinate covers the cell; with
one or two active coordinates it is absorbed into a unary mask or pairwise
constraint. Boxes with at least three active coordinates remain hard.
The following is the essential control flow. At the root, \texttt{remaining}
is unset. Each recursive step advances the axis and consumes one level of
the current block, even when no cut is available.

\begin{quote}\small
\begin{verbatim}
Node(boxes, state, cell, axis, remaining):
    if state is empty: return 0
    hard, state = absorb_easy(boxes, state, cell)
    if state is empty: return 0
    if len(hard) <= 2: return base_sum(...)
    if remaining == 0:
        state, hard, cell = compress_to_blocks(...)
        remaining = unset
    if remaining is unset:
        remaining = restart_block_length(...)
    cut = choose_weighted_median(hard, cell, axis)
    next_axis, next_budget = (axis + 1) % 4, remaining - 1
    if cut is absent:
        return Node(hard, state, cell, next_axis, next_budget)
    left, right = split(cell, axis, cut)
    return Node(hard, state, left, next_axis, next_budget)
         + Node(hard, state, right, next_axis, next_budget)
\end{verbatim}
\end{quote}

The returned quantity is uncovered mass; the public computation subtracts it
from the total mass of the enclosing product grid. At a base cell, at most
two hard boxes remain. Their inclusion--exclusion therefore needs at most
four rectangle contractions. The complete small-base loop is compiled.
The configured base threshold can be changed, but a compiled base with 63
or more hard boxes is rejected to prevent native bitmask overflow.

Compression groups fine cells into consecutive blocks. In 4D it lifts the
state to eight variables---four fine indices and four block indices---and
imposes the inclusive membership bounds. Eliminating the four fine indices
leaves four block indices whose weights represent sums over the old cells.
The same mass-summing operation handles signed intermediate states and the
magnitude atom. The existing potential-based block schedule, weighted median
rule, axis cycling and base threshold are retained.

The subclass introduces only four numerical hooks:
\begin{center}\small
\begin{tabular}{lp{98mm}}\toprule
Hook & Responsibility\\\midrule
\texttt{\_initial\_terms} & Pack the initial state once.\\
\texttt{\_apply\_easy} & Apply all classified unary and pairwise updates in Numba.\\
\texttt{\_base} & Evaluate the complete small-base inclusion--exclusion in Numba.\\
\texttt{\_compress} & Sum fine indices into blocks in Numba and retain native output.\\
\bottomrule\end{tabular}
\end{center}
Python performs geometric decisions and constructs small mask and block
descriptions. It does not unpack the term collection after compression or
repack it at every leaf. The geometric recursion is shared with the hybrid
backend rather than copied into a second implementation.

\section{Why the native state matters}
The first compiled version still packed Python term objects into native
containers at each base contraction and converted compressed output back
into Python objects. Its warm improvement over the hybrid was only
1.12--2.19 times across the five matched cases, with median 1.59.
A separate profile of the seed-42 twenty-point sphere exposed the remaining cost:

\begin{table}[ht]\centering\small
\begin{tabular}{lrr}\toprule
Profile statistic & Repeated conversion & Persistent native state\\\midrule
Total profiled seconds & 2.819 & 0.265\\
Python-visible function calls & 4,129,144 & 28,952\\
\texttt{pack()} calls & 41 & 1\\
\texttt{pack()} cumulative seconds & 2.306 & 0.001027\\\bottomrule
\end{tabular}
\caption{Independent warm profiles. Native work is not counted as individual
Python calls; these counts do not measure arithmetic complexity.}
\label{tab:profile}
\end{table}

Packing accounted for approximately 81.8\% of the first profile. Keeping
native term states through the recursion removed the repeated conversion.
The final profile spends about 0.175 seconds in base contractions, 0.045
seconds in compression, and 0.030 seconds applying easy constraints.
These profiled durations include instrumentation overhead and are not used
as benchmark speedups. The change is supported by both the profile and the
subsequent independent timing experiment.

\section{The 4D timing experiment}
Table~\ref{tab:timings} reports all seven input sets. Each algorithm receives
the same saved points. New and hybrid backends receive anchored extents $r-p$;
the older prefix and Python Chan implementations receive positive minimization
coordinates and reference $r$. The conversion is included in the new engines'
timed calls. No dominance prefilter is applied, and the base threshold is two
where configurable.

Workers run sequentially, with one fresh process per algorithm/input pair.
After an untimed first call, a fresh solver is created for each of three
timed calls; their median is reported. Timed work includes conversion,
preprocessing, packing and evaluation. Compilation caches were populated by
correctness checks beforehand. Each worker has a 45-second total budget,
including imports, first use and all repetitions.

\begin{table}[ht]\centering\small
\setlength{\tabcolsep}{4pt}
\begin{tabular}{lrrrrr}\toprule
Family seed / $n$ & New JIT & Hybrid & Old prefix & Chan $d/3$ & Chan $d/2$\\\midrule
Sphere 42 / 10 & 4.924 & 12.331 & 6.478 & 9.437 & 1.408\\
Sphere 42 / 20 & 252.106 & 2050.424 & 19.284 & 34.166 & 5.145\\
Sphere 42 / 40 & 7540.933 & budget & 55.884 & 134.478 & 18.260\\
Sphere 99 / 20 & 274.465 & 2172.687 & 21.673 & 32.210 & 5.612\\
Sphere 99 / 40 & 6929.245 & budget & 51.891 & 122.594 & 17.727\\
Grid 7 / 20 & 3.414 & 5.013 & 0.766 & 2.263 & 0.362\\
Uniform 17 / 20 & 115.610 & 897.588 & 14.112 & 18.851 & 2.529\\
\bottomrule\end{tabular}
\caption{Warm end-to-end medians in milliseconds. ``budget'' is a whole-worker
timeout; it is not a per-call lower bound. Both Chan reference codes use Python
without Numba.}
\label{tab:timings}
\end{table}

The five completed matched comparisons yield speedups of
1.47--8.13, with median ratio 7.76. For the seed-42 twenty-point sphere,
2.050 seconds becomes 0.252 seconds, an 8.13-fold improvement. Both forty-point
spheres complete in the new backend, at 7.541 and 6.929 seconds, whereas the
hybrid workers exhaust the budget. No ratio is inferred from those timeouts.

There are 33 completed algorithm/input combinations out of 35, with no
numerical mismatches against the independent Python Chan $d/2$ evaluator.
Maximum scaled disagreement is $7.89\times10^{-16}$, using
$|x-y|/\max(1,|y|)$. All seven new-backend cases exercised nopython
base evaluation and easy updates; the five cases requiring compression also
exercised nopython block compression. Fastmath and parallel execution are disabled.

\textbf{The older compiled prefix solver remains substantially faster on the
larger spherical inputs.} It takes about 19 milliseconds for seed-42 $n=20$
and 56 milliseconds for $n=40$. Python Chan $d/2$ is fastest on all seven
inputs. Three warm repetitions on small inputs neither determine asymptotic
growth nor establish an eventual crossover; no exponent is fitted here.

The machine is an AMD Ryzen 7 5800H with eight physical cores, sixteen logical
processors and approximately 16 GiB RAM, running Windows 11. Software versions
are CPython 3.12.14, NumPy 2.5.3 and Numba 0.68.0. Sphere points are absolute
Gaussian vectors normalized to unit Euclidean length, with $r=(1.2,1.2,1.2,1.2)$.
Seed 42 consumes sizes $5,10,20,40$ successively, discarding the five-point
case; seed 99 consumes $20,40$. The grid case uses independent integer
coordinates in $\{0,1,2,3\}$ divided by four, and the uniform case uses
independent coordinates in $[0,1)$; both have $r=(1,1,1,1)$.
All actual inputs are archived. No confidence intervals or processor pinning
were used.

\subsection*{First use and compilation}
A separate process with an empty dedicated Numba cache measured startup on
the seed-42 twenty-point sphere. Imports took 0.803 seconds. The first
computation took 48.408 seconds, including compilation and evaluation; the
next call took 0.211 seconds. Values and counters agreed, and all three
compiled entry points had nopython signatures. These are single observations,
not medians, and the first duration is not compiler time alone.

The warm speedup therefore has a substantial first-use cost. Persistent
caching is enabled. The cold-start script requires a new or empty cache
directory and does not delete existing caches. Its input, source hashes and
full measurements are archived separately from the warm benchmark.

\section{Correctness and direct magnitude support}
The dedicated verification suite checks the numerical kernels against their
defining finite sums and compares small complete problems with exact rational
inclusion--exclusion. Its main coverage is:
\begin{center}\small
\begin{tabular}{lr}\toprule
Check & Count\\\midrule
Exact small HV/magnitude instances & 50\\
Signed rectangle contractions & 48\\
Signed easy-constraint cells & 16\\
First-compression cells & 144\\
Masked repeated-compression checks & 16\\
Default twenty-point hybrid comparisons & 2\\
Exact twelve-point cases reaching generation two & 2\\
Affine magnitude/HV comparisons & 9\\\bottomrule
\end{tabular}
\end{center}
These counts describe different assertions and are not a count of independent
large end-to-end instances. Further checks cover empty sets, duplicates, zero
coordinates, tied inputs, mixed scales, input immutability, term cancellation,
hash collisions, the collision-search cap, native bitmask overflow, public
API compatibility, and a fresh-process package import using cached types.
All checks pass. Maximum recorded scaled error in the full suite is
$2.68\times10^{-14}$; this is empirical evidence rather than a universal
floating-point error bound.

The native state is packed exactly once in each of the 52 nonempty complete
computations instrumented for this invariant, and zero times for the two empty
ones. The twenty-point checks perform ten compressions each; the forced
twelve-point checks reach a second compression generation. Thus successful
tests are not restricted to cases in which compression is bypassed.

Magnitude is also checked through the independent affine identity
\begin{equation}
 \operatorname{Mag}(D(A))=
 \operatorname{HV}_4\!\left(\{(1+a_1/2,\ldots,1+a_4/2):a\in A\}\right).
 \label{eq:mag}
\end{equation}
Every nonempty collection $S\subseteq A$ of anchored boxes has intersection
with upper corner $b_i=\min_{a\in S}a_i$. Its product-measure value is
$\prod_i(1+b_i/2)$, which is the ordinary volume of the transformed
intersection. Inclusion--exclusion proves~\eqref{eq:mag}. This identity
concerns anchored-set values; it does not assert translation invariance of
the atomic measure.

The reverse direction is equally explicit. If every coordinate of every
corner in $C$ is at least one, set $a_i=2(c_i-1)$. Then
\[
 \operatorname{HV}_4(C)=\operatorname{Mag}(D(2(C-\mathbf 1))).
\]
For general nonnegative corners, first discard boxes with a zero extent:
they contribute no four-dimensional volume. If no boxes remain, the answer
is zero. Otherwise, choose a positive scale $s$ so that every coordinate of
every corner in $sC$ is at least one. Homogeneity gives
\[
 \operatorname{HV}_4(C)=s^{-4}
 \operatorname{Mag}\bigl(D(\{2(sc-\mathbf 1):c\in C\})\bigr).
\]
The preprocessing takes linear work. These identities preserve the
underlying exact-value computation up to transformations; they do not remove
the signed-term expansion or guarantee equal floating-point conditioning.

For the four generators
\[
 \{(1,4,3,2),(2,3,4,1),(3,2,1,4),(4,1,2,3)\},
\]
the public interface returns ordinary HV $69$ and magnitude $765/16$.
The timing table concerns ordinary HV; it does not supply magnitude timing
ratios merely because both modes pass correctness checks.

\section{What still limits performance}
Compilation leaves the number of emitted terms unchanged on the measured
matched cases. On seed-42 $n=20$, both hybrid and new versions use 61 nodes,
ten compressions, 10,809 raw terms, 7,441 term eliminations and a peak of 625
terms. The new $n=40$ run uses 157 nodes, thirteen compressions, 251,360 raw
terms, 146,333 eliminations and a peak of 7,501 terms.

The older prefix solver uses a different representation: six staircases and
a compiled target-plane scan. It does not perform this scheduled signed-term
compression. Its smaller measured cost is consequently not a contradiction
of the new backend's successful compilation. Nor does its speed establish
the complete $\widetilde O(n^{4/3})$ bound for that implementation; a global
bound on its target-grid work remains separate from benchmarking.

The useful implementation result is precise: complete contractions can run
in native code while the recursion stays short and readable, and retaining
native state is essential to realizing the benefit. The remaining target
is to reduce or better control the states produced by compression. Changing
the schedule, reusing equivalent summaries, or introducing a specialized
planar contraction would require its own correctness and complexity analysis.
None of those additional improvements is claimed in this report.

\section{Use and reproduction}
The established 4D Numba entry point now selects the new backend:
\begin{quote}\small
\begin{verbatim}
from numerical_chan4_numba import hypervolume4
value = hypervolume4(corners)
mag = hypervolume4(corners, magnitude=True)
\end{verbatim}
\end{quote}
Keep the \texttt{NumericalChanHV4D} and \texttt{NumericalChanHVND} directories
together. The adapter is \texttt{numerical\_chan4\_compiled.py}; the compiled
algebra is \texttt{compiled\_terms.py}. Both script and package imports use
one canonical kernel module name so cached native tuple types remain importable.
Exact arithmetic remains available through \texttt{numerical\_chan4.py}
with suitable integer or rational inputs.

From the repository root, install the requirements in the 4D directory and run
\texttt{python NumericalChanHV4D/verify\_whole\_loop.py}. This also populates
the cache before warm benchmarking. The benchmark directory
\texttt{NumericalChanHV4D/benchmarks/whole\_loop} contains:
\begin{itemize}
\item \texttt{benchmark\_whole\_loop.py}, accepting \texttt{--old} with the
path to the pinned older \texttt{HV4DMagnitude} checkout;
\item \texttt{profile\_whole\_loop.py}, for a separate warm profile;
\item \texttt{cold\_start.py}, accepting \texttt{--cache-dir} with a new or
empty directory;
\item the saved inputs, repetitions, counters, source hashes, profiles and
startup record in \texttt{results}, and the first compiled version in
\texttt{results\_v1}.
\end{itemize}

The measured native-state sources are pinned at \texttt{1c2b575}; the first
compiled version is \texttt{7ea05d2}~\cite{timed}. Release \texttt{0c4375e} also selects
the new public wrapper and fixes the canonical import path. That import-only
change leaves the measured numerical kernel unchanged. The cold-start record
uses the release import convention. The independent Chan reference files and
older prefix sources are identified by their archived hashes and commits.
Full links appear below; this report itself does not require an external
BibTeX database or additional figures to compile.

\begin{thebibliography}{9}\raggedright
\bibitem{chan} T. M. Chan. \emph{Klee's Measure Problem Made Easy}.
FOCS 2013, pp.~410--419. DOI: 10.1109/FOCS.2013.51.
\href{https://tmc.web.engr.illinois.edu/easyklee8_13.pdf}{Author's August 2013 version},
especially Sections 4.1--4.2.
\bibitem{magnitude} M. T. M. Emmerich.
\emph{The Magnitude of Dominated Sets: A Pareto Compliant Indicator Grounded
in Metric Geometry}. arXiv:2604.18147, 2026.
\url{https://arxiv.org/abs/2604.18147}.
\bibitem{old} M. Emmerich and collaborators. \emph{HV4DMagnitude},
source snapshot \texttt{0c668e765202}.
\href{https://github.com/emmerichmtm/HV4DMagnitude/tree/0c668e76520201f41ac0f144caa8d468d6c84e86}{Pinned repository}.
The measured prefix path uses \texttt{solver.py}, \texttt{prefix4.py} and
\texttt{fastkernel.py}.
\bibitem{release} \emph{FastHVChan / NumericalChanHV4D}, release
\texttt{0c4375e963c8}.
\href{https://github.com/emmerichmtm/FastHVChan/tree/0c4375e963c8c054b7ebba10b715ba55de6e9529/NumericalChanHV4D}{Code, tests, implementation guide and benchmark data}.
The repository's Python reference codes are implementations of Chan's
approaches; they are not identified here as code written by Chan.
\bibitem{timed} \emph{FastHVChan}, numerical source snapshots:
\href{https://github.com/emmerichmtm/FastHVChan/tree/1c2b5756477e4bc428b0b49acd4b7e9cc1526609}{native-state version \texttt{1c2b575}};
\href{https://github.com/emmerichmtm/FastHVChan/tree/7ea05d21b67de869268740096fad32a1117b0430}{first compiled version \texttt{7ea05d2}}.
Raw data and SHA256 hashes are archived under \texttt{benchmarks/whole\_loop}
in release~\cite{release}.
\end{thebibliography}
\end{document}
